\problem{Q01}{智能行为、演绎与归纳：怎样写概念短答}{正面第1栏 F01；相关新课件实际第7--12、25、29--35页}
\qtext{以下均为自拟复习题。一个分拣系统有两种工作方式：甲按预设规则“损伤超过阈值则剔除”处理水果；乙根据带标签的水果样本调整模型，再判断新水果。\begin{enumerate}
\item 说明智能系统的含义，并区分甲所用的演绎与乙所用的归纳。
\item 比较 Artificial Intelligence (AI) 与 Computational Intelligence (CI)，以及 hard computing 与 soft computing。评价“AI 不学习，只有 CI 才学习”是否准确。
\item 将 McCulloch--Pitts、Rosenblatt、Widrow、Werbos 与本课模型／方法联系起来；只需人物--贡献，不要求争议年代。
\end{enumerate}}
\route{圈出“预设规则”“从样本”“新水果”“AI/CI”“区别”。查 F01 的定义和比较；本题没有数值目标或更新规则，不应套感知机、BP 或最优化计算式。}
\given{规则与当前事实是甲的已知；样本及其标签是乙的经验；新水果是未见输入。这里比较的是获得／使用知识的方式，不是比较水果特征的数值。}
\solve
\begin{enumerate}
\item 智能系统是表现出某些智能行为的系统，例如推理、从经验学习、适应环境、规划或设计；不要求同时具备全部人类能力。甲从规则和当前事实推出处理结论，是\kw{演绎 deduction}。乙从有限实例形成可用于其他实例的关系，是\kw{归纳 induction}，其推广结论仍可能出错。学习后利用固定参数算新输入，另称预测／推理；此处 inference 不必是形式逻辑证明。
\item AI 是研究和实现智能能力的广泛领域。本课 CI 重点讨论 neural networks、fuzzy logic、evolutionary computation，分别强调学习、对不精确信息的表示及优化等互补能力。课件的传统 AI--CI 对照强调符号／显式规则与适应性取向，但它们并非现代意义下永不交叠的集合；神经网络学习也属于广义 AI。
\item Hard computing 的典型取向是明确模型、精确条件与清晰规则；soft computing 容纳不确定性、不精确性和近似解。“软”不表示没有严格数学，也不保证任何软计算方法都优于精确方法。应依任务、数据和资源选择。
\end{enumerate}
\begin{center}\begin{tabular}{ll}\toprule
人物 & 本课联系\\\midrule
McCulloch--Pitts & 二值神经元与阈值逻辑\\
Frank Rosenblatt & 感知机 Perceptron\\
Bernard Widrow & ADALINE 与误差最小化学习\\
Paul Werbos & 反向传播相关早期工作\\\bottomrule
\end{tabular}\end{center}
\ans{演绎依据给定规则推出结论；归纳从实例建立可推广关系，但需新数据检验。CI 是本课强调神经、模糊、进化方法的家族，能够与其他 AI 方法结合。AI 与学习并不互斥；硬／软计算描述方法取向，不能用来代替具体任务评价。}
\qcheck{逐句检查：是否分别写清“规则\(\to\)结论”和“样本\(\to\)一般关系”；是否提到归纳的不确定性；是否避免把 AI/CI 或 hard/soft 写成绝对互斥。人物表应与模型名称对应。}
\pitfall{用“都是人工智能”代替明确比较；把一个下棋系统的成功写成通用智能证明；把导论中出现的模糊、进化或搜索名称当成本次算法计算范围。}
\transfer{若乙只把原样本存成查找表，要另检验其对未见输入的表现。若系统同时执行规则并训练网络，两种工作方式可以共存，不必二选一。}

\problem{Q02}{ANN 统一框架：结构、五条规则与分类轴}{正面第1栏 F02；正面第2栏 N02；新课件实际第36--50页}
\qtext{某网络有 \(d=2\) 个外部输入、\(N=3\) 个状态节点、\(M=2\) 条可调连接和 \(m=1\) 个外部输出。两个输入节点传值，第三节点聚合加权输入，再经确定性非线性函数；输出选第三节点状态。训练提供目标标签，逐样本调整权重；预测时权重固定，运行连接无环。解释 \(\mathrm{ANN}=(\mathrm{ARCH},\mathrm{RULE})\)，列出有关变量及五条规则，并按课件四个分类轴判断该网络。最后说明训练、推理、状态变化、权重学习的区别。}
\route{圈“变量／规则”“有标签”“确定性”“无环”“固定权重”。查 F02；若问反馈与状态再看 N02。不能由“非线性”推成循环，也不能由“逐样本更新”推成随机激活。}
\given{\(\bm x\in\R^2\) 是外部输入，\(\bm u,\bm v\in\R^3\) 是净输入和激活状态，\(\bm w\in\R^2\) 收集两条连接权重，\(y\in\R\) 是读出。\(z\) 为训练目标，不是预测时必需的输入。}
\solve
ARCH 指部件、连接以及输入输出的组织；RULE 指它们如何运算和学习。只抄五个结构字母仍不足以复现网络，还要说明连接方向、权重及函数。
\begin{center}\begin{tabular}{p{.12\linewidth}p{.77\linewidth}}\toprule
规则 & 本题可对应的实际操作\\\midrule
\(G\) 聚合 & 从输入／旧状态与权重形成净输入，如 \(u_3=w_1x_1+w_2x_2\)。\\
\(F\) 激活 & \(v_3=f(u_3)\)，输入节点直接传值。\\
\(H\) 读出 & \(y=v_3\)，从内部状态取得外部输出。\\
\(E\) 评价 & 例如在给定标签时计算 \(E=\tfrac12(z-y)^2\)。\\
\(L\) 学习 & 根据具体学习规则产生 \(\Delta\bm w\)，令 \(\bm w_{k+1}=\bm w_k+\Delta\bm w_k\)。\\\bottomrule
\end{tabular}\end{center}
实际运行顺序通常是 \(G\to F\to H\)；训练时再评价并更新。\(\mathrm{RULE}=(E,F,G,H,L)\) 是组成清单，非必须按字母顺序执行。

该网络在四轴上是：\kw{确定性}（固定条件给出同一激活）、\kw{前馈}（运行连接无环）、\kw{半线性聚合}（加权和再激活，整体仍可非线性）、\kw{监督学习}（有目标标签）。若聚合含 \(x_1x_2\) 等交互项，才属于本课所说高阶聚合；“隐藏层多”是另一维度。

训练利用数据改变参数；推理以现有参数计算输出。输入变化会使 \(\bm u,\bm v,y\) 变化，并不自动改变 \(\bm w\)。课件的快状态／慢权重对应两种时间尺度，不等于任何前馈网络都有历史记忆。
\ans{ANN 由结构与规则共同定义；本题变量维度分别为 \(2,3,3,2,1\)（按 \(x,u,v,w,y\)）。五规则为聚合、激活、读出、评价和学习。本题是确定性、前馈、加权和型半线性聚合、监督网络；预测中的状态变化与训练中的权重变化是不同动作。}
\qcheck{\(G\) 输出应与 \(u\) 同维，\(F\) 输出与 \(v\) 同维，\(H\) 输出与 \(y\) 同维；权重更新前后维度必须一致。四轴分别检查，不把四个名称当四选一。}
\pitfall{把 \(u\) 与 \(v\) 永远当同一量；以为课件 \(E,L\) 简写没列标签就不需要目标；把实数输入等同连续时间，或把误差用于训练的反馈箭头当作运行连接的循环。}
\transfer{若只去掉标签，监督性必须重判；若只加入状态反馈，前馈／循环必须重判。改变权重初始化方式不必改变前向激活是否确定。连续时间式 \(d\bm w/dt=L\) 给的是变化率，小步更新还需乘时间步长。}

\problem{Q03}{从任务选择输出、监督信息与模型层次}{正面第1栏 F01/F02；正面第2栏 N02；模型细节另见相应栏}
\qtext{对下列任务分别写出需要的输出、是否提供监督信息，以及应先考虑什么模型性质。A：给带“合格／不合格”标签的水果学习分类；B：给实测下一小时负荷学习预测；C：无标签地整理相似观测；D：相同输入脉冲下，输出须取决于先前状态；E：规定输入三个二值开关恰为 \((1,0,1)\) 时报警。评价“只要训练误差最低，所有任务都选同一模型”及“BP 与 MLP 是两种并列网络”两句话。无需替范围外算法设计训练过程。}
\route{圈“标签”“实测数值”“无标签”“先前状态”“规定模式”。先用 F01/F02 区分任务和信息，再查模型。分类／回归由目标输出决定，不能仅凭题中出现“预测”二字判断。}
\given{A 的目标为类别；B 的 \(y\) 是连续实测量；C 只有输入 \(\bm x\)，没有本任务正确类别；D 还需内部状态 \(s_t\)；E 的逻辑要求完全给定。没有给定真实数据分布，不能声称唯一最优模型。}
\solve
\begin{enumerate}
\item A 是监督分类，需要类别或判别分数再决策。若原特征线性可分，可先考虑感知机／线性分类；边界复杂时应考虑非线性表示。需用留出数据评价推广表现。
\item B 是监督回归，输出实数。可比较范围内 RBF、SVR 或 LS-SVM 等回归形式；选择还取决于损失、正则化、数据规模及验证结果。“预测”不自动等于给类别。
\item C 是无监督的结构发现任务，输出可以是组别或简化表示。无监督仍需相似性标准和评价；组号不自动对应“合格类别”。本题只识别任务，不展开后续聚类算法。
\item D 需要历史信息。可用状态递推 \(s_{t+1}=T(s_t,x_t)\) 描述循环结构，必须规定初态和更新方式；普通无状态前馈输入若不含历史，无法区分相同当前输入的不同过去。若显式把历史窗口加入输入，也可用前馈映射处理，需说清输入已改变。
\item E 适合直接构造 M-P／阈值逻辑 decoder。因为目标规则已明确，不必先收集样本让网络学习同一个真值表；N01 给出可核对构造。
\end{enumerate}
模型选择至少要分清三个问题：\kw{表示能力}决定模型能否表达所需关系；\kw{优化}决定训练能否找到较好参数；\kw{泛化}决定未见数据上能否有效。训练误差最低不足以回答后两者之间的权衡，也可能反映过拟合。

MLP 是多层感知机模型／结构；BP 是计算其参数梯度的反向传播方法，随后可由梯度下降等规则更新参数。二者所在层次不同，不能当同类模型二选一。
\ans{A：监督分类；B：监督回归；C：无监督结构发现；D：带状态的映射或显式历史特征；E：已知规则的逻辑构造。选模型应同时检查任务、输入信息、表示、损失与验证表现。BP 是训练中的梯度计算方法，MLP 是模型。}
\qcheck{把每项写成“输入信息\(\to\)输出对象”，检查输出究竟是类别、实数、组别、状态还是逻辑值。对任何推荐加上所需条件，不能把某名称写成无条件最佳。}
\pitfall{“无标签＝没有目标函数”“有反馈＝正在学习”“回归输出再套分类阈值”“MLP与BP并列比较”。范围内的模型选择题可以提到无监督用途，但无需补入本课件未展开的后续算法。}
\transfer{若 B 改成预测“高负荷／低负荷”，应重定义标签与损失；若 A 的输入已含过去十步观测，不能仅凭时间数据就断言模型必须有循环连接。}

\problem{Q04}{向量维度、加权和与仿射映射}{正面第1栏 M01/F02；正面第2栏 N01}
\qtext{给定列向量与参数
\[\bm x=\begin{bmatrix}2\\-1\end{bmatrix},\quad
W=\begin{bmatrix}1&2\\-1&3\end{bmatrix},\quad
\bm b=\begin{bmatrix}1\\-2\end{bmatrix}.
\]
先算 \(\bm u=W\bm x+\bm b\)，再逐分量用 \(H(a)=1\)（\(a\ge0\)）、否则0，得 \(\bm v\)。读出 \(y=[1\ 0]\bm v\)，训练目标 \(z=0\)，评价 \(E=\tfrac12(z-y)^2\)。写出各量维度和完整前向计算。说明 \(W\bm x\)、\(W\bm x+\bm b\)、整个带阈值网络各属于线性、仿射还是非线性。}
\route{圈“列向量”“矩阵乘法”“bias”“逐分量激活”“读出”。查 M01 的维度和仿射定义、F02 的 G/F/H/E；阈值等号查 N01。题目未要求更新权重，不使用学习规则。}
\given{\(\bm x,\bm b,\bm u,\bm v\in\R^2\)，\(W\in\R^{2\times2}\)，读出行向量在 \(\R^{1\times2}\)，\(y,z,E\in\R\)。矩阵的每一行与同一个输入作内积。}
\solve
\begin{enumerate}
\item 聚合：\((2\times2)(2\times1)\) 得 \(2\times1\)，可与同形状 \(\bm b\) 相加。
\[
W\bm x=\begin{bmatrix}1(2)+2(-1)\\-1(2)+3(-1)\end{bmatrix}
=\begin{bmatrix}0\\-5\end{bmatrix},\qquad
\bm u=\begin{bmatrix}0+1\\-5-2\end{bmatrix}=\begin{bmatrix}1\\-7\end{bmatrix}.
\]
\item 激活：\(u_1=1\ge0\)，故 \(v_1=1\)；\(u_2=-7<0\)，故 \(v_2=0\)。于是 \(\bm v=(1,0)\T\)。
\item 读出与评价：\(y=1v_1+0v_2=1\)，\(e=z-y=-1\)，所以 \(E=\tfrac12(-1)^2=\tfrac12\)。误差 \(e\) 与非负代价 \(E\) 是不同量。
\item \(\bm x\mapsto W\bm x\) 为线性变换；加 \(\bm b\ne0\) 后为仿射变换。仿射映射不必过原点；本题在 \(\bm x=\bm0\) 时输出净输入 \(\bm b\ne0\)，已不符合线性变换必须保持零向量的条件。
\item 阶跃激活为非线性，因此整个输入到输出映射一般非线性；不能因聚合是加权和就断言整个网络线性。
\end{enumerate}
\ans{\(W\bm x=(0,-5)\T\)，\(\bm u=(1,-7)\T\)，\(\bm v=(1,0)\T\)，\(y=1,E=1/2\)。前两种映射分别为线性与仿射，加入阈值激活后的整体映射一般非线性。}
\qcheck{用行内积独立重算两个净输入；检查“行数决定输出维数”。所有比较在加完偏置之后执行。若误把 \(W\T\bm x\) 当本题聚合，将得到不同结果，不能只因形状恰好相容就交换。}
\pitfall{把矩阵乘法当逐元素相乘；漏加负偏置；把 \(z-y\) 与其平方误差混同；把非零偏置的映射直接称为线性。这里一般负权按乘法处理，没有 M-P veto。}
\transfer{若题设阈值为 \(\theta_j\)，先比较 \(u_j-\theta_j\)，或把它并入偏置；若题设输出为净输入而没有激活，不应擅自加阶跃。加偏置可写增广输入 \(\tilde{\bm x}=(\bm x\T,1)\T\)。}

\problem{Q05}{链式法则、平方误差梯度与驻点性质}{正面第1栏 M02/M01}
\qtext{设
\[A=\begin{bmatrix}1&2\\0&1\end{bmatrix},\quad
\bm z=\begin{bmatrix}1\\-1\end{bmatrix},\quad
J(\bm w)=\tfrac12\norm{A\bm w-\bm z}^{2},\quad \bm w\in\R^2.
\]
用链式法则求梯度；计算 \(\bm w=(0,0)\T\) 时的梯度；求驻点并判断是最大、最小还是不能判断。另写出 \(\tfrac12\bm w\T Q\bm w\) 的一般梯度，并评价“所有梯度为零的点都是最小值”。}
\route{圈“平方误差”“求导”“驻点”“最大／最小”“二次型”。查 M02，先确定目标方向及变量维度。本题是目标极值判断的复习设计，不推断“最高／最大”具体指什么，也不把所有模型都当最大化。}
\given{令 \(\bm r=A\bm w-\bm z\in\R^2\) 为预测减目标的残差；\(A\in\R^{2\times2}\)。本题 \(J=\tfrac12(r_1^2+r_2^2)\)，前面的 \(1/2\) 会抵消平方求导产生的2。}
\solve
把复合计算拆成 \(\bm w\to\bm r\to J\)。\(\partial J/\partial r_i=r_i\)，而 \(\partial r_i/\partial w_j=A_{ij}\)。对影响同一个 \(w_j\) 的全部残差贡献求和：
\[
\frac{\partial J}{\partial w_j}=\sum_i A_{ij}r_i,
\qquad \nabla J=A\T(A\bm w-\bm z).
\]
明确展开，可独立检查：
\[
J=\tfrac12(w_1+2w_2-1)^2+\tfrac12(w_2+1)^2,
\quad \nabla J=\begin{bmatrix}w_1+2w_2-1\\2w_1+5w_2-1\end{bmatrix}.
\]
在 \(\bm w=\bm0\)，\(\bm r=(-1,1)\T\)，梯度为 \((-1,-1)\T\)。驻点满足
\[
w_1+2w_2=1,\quad 2w_1+5w_2=1.
\]
第二式减第一式的2倍得 \(w_2=-1\)，再回代得 \(w_1=3\)。此时 \(A\bm w=\bm z\)，所以 \(J=0\)。由于平方和非负，零值已经证明它是全局最小值；且 \(A\) 可逆，残差全零解唯一。

二次型需注意矩阵对称性：
\[
\nabla_{\bm w}\!\left(\tfrac12\bm w\T Q\bm w\right)
=\tfrac12(Q+Q\T)\bm w.
\]
仅当 \(Q=Q\T\) 时才能简化为 \(Q\bm w\)。本题 Hessian 为
\(A\T A=\begin{bmatrix}1&2\\2&5\end{bmatrix}\)，对非零 \(\bm h\)，\(\bm h\T A\T A\bm h=\norm{A\bm h}^2>0\)，也验证严格凸性。
\ans{\(\nabla J=A\T(A\bm w-\bm z)\)；原点梯度为 \((-1,-1)\T\)；唯一驻点 \((3,-1)\T\) 是全局最小点，最小值0。驻点本身不自动说明极值性质：\(t^2\) 在0最小，\(-t^2\) 在0最大，\(t^3\) 在0无局部极值。}
\qcheck{将 \((3,-1)\T\) 代回两条梯度式以及 \(A\bm w=\bm z\)，均应成立。维度 \((2\times2)(2\times1)=2\times1\)；若写成 \(A\bm r\)，形状在此例虽相容，链式方向仍错。}
\pitfall{把残差正负与下降方向混淆；漏掉转置；不检查 \(Q\) 对称就写 \(Q\bm w\)；只由 \(\nabla J=0\) 就宣称全局最优。}
\transfer{若损失去掉 \(1/2\)，梯度整体乘2；若加 \(\lambda\norm{\bm w}^2/2\)，梯度加 \(\lambda\bm w\)。若 \(A\) 列不独立，平方和仍为凸目标，但最小点未必唯一。}

\problem{Q06}{增广矩阵、消元与三种解的状态}{正面第1栏 M03/M01}
\qtext{分别把以下系统写成 \(A\bm q=\bm c\)，用增广矩阵行变换求解，并判断唯一解、无解或无穷多解；不使用 \(A^{-1}\bm c\) 作为唯一解法。
\[
\mathrm{I}:\ \begin{cases}2x+y=5,\\x-y=1;\end{cases}\qquad
\mathrm{II}:\ \begin{cases}x+y=2,\\2x+2y=5;\end{cases}\qquad
\mathrm{III}:\ \begin{cases}x+y=2,\\2x+2y=4.\end{cases}
\]}
\route{圈“多元一次”“矩阵”“消元”“判断解”。查 M03：先按同一未知量顺序排列列，再做不改变解集的行操作；看到方程数等于未知数，仍须检查独立性与一致性。}
\given{三个系统均取 \(\bm q=(x,y)\T\in\R^2\)。系数按 \(x,y\) 两列排列，最后一列是右端。\(R_i\) 表示第 \(i\) 行，增广分隔线不代表额外未知量。}
\solve
\kw{I：}\(A=\begin{bmatrix}2&1\\1&-1\end{bmatrix}\)，\(\bm c=(5,1)\T\)。
\[
\left[\begin{array}{cc|c}2&1&5\\1&-1&1\end{array}\right]
\xrightarrow{R_1\leftrightarrow R_2}
\left[\begin{array}{cc|c}1&-1&1\\2&1&5\end{array}\right]
\xrightarrow{R_2\leftarrow R_2-2R_1}
\left[\begin{array}{cc|c}1&-1&1\\0&3&3\end{array}\right].
\]
第二行给 \(3y=3\)，所以 \(y=1\)；回代 \(x-y=1\) 得 \(x=2\)。两列均有主元，\(\operatorname{rank}A=\operatorname{rank}[A|\bm c]=2\)，故唯一解。

\kw{II：}\(A=\begin{bmatrix}1&1\\2&2\end{bmatrix}\)，\(\bm c=(2,5)\T\)。
\[
\left[\begin{array}{cc|c}1&1&2\\2&2&5\end{array}\right]
\xrightarrow{R_2\leftarrow R_2-2R_1}
\left[\begin{array}{cc|c}1&1&2\\0&0&1\end{array}\right].
\]
第二行要求 \(0x+0y=1\)，矛盾。\(\operatorname{rank}A=1\)，增广矩阵秩2，故无解；不应继续把0当除数。

\kw{III：}\(A\) 与II相同，\(\bm c=(2,4)\T\)。
\[
\left[\begin{array}{cc|c}1&1&2\\2&2&4\end{array}\right]
\xrightarrow{R_2\leftarrow R_2-2R_1}
\left[\begin{array}{cc|c}1&1&2\\0&0&0\end{array}\right].
\]
第二行是恒等式，只有一个独立约束。令自由变量 \(y=t\)，则 \(x=2-t\)，\(t\in\R\)。两矩阵秩同为1，小于未知数2，故无穷多解。
\ans{I：\((x,y)=(2,1)\)，唯一解。II：出现矛盾行 \([0\ 0|1]\)，无解。III：\((x,y)=(2-t,t)\)，\(t\in\R\)，无穷多解。一般先查一致性；一致且主元数等于未知数则唯一，一致但有自由变量则无穷多。}
\qcheck{I 回代得 \(2(2)+1=5\)、\(2-1=1\)；III 回代得 \((2-t)+t=2\) 及其2倍4。II 原式第二左端必须是第一左端2倍，右端却不是，直接暴露矛盾。}
\pitfall{只改系数不改右端；把 \([0\ 0|0]\) 与 \([0\ 0|1]\) 混同；认为方程数等于未知数就一定唯一；列交换后忘记同步交换未知量解释。}
\transfer{把II右端5改成4，矛盾消失但没有增加独立约束，因此变成无穷多解。后面 LS-SVM 的分块系统同样按已声明未知量列顺序消元和回代。}

\problem{Q07}{M-P 的 veto、逻辑门与阈值边界}{正面第2栏 N01；新课件实际第51--56页}
\qtext{所有输入均为 \(\{0,1\}\)，达到阈值输出1。\begin{enumerate}
\item M-P 单元的前两个输入兴奋、第三个输入抑制，阈值 \(\theta=1\)。求 \(\bm x=(1,1,1)\) 的输出；再与普通加权 TLU 的 \(\bm w=(1,1,-1)\)、\(\theta=1\) 比较。
\item 用 M-P 单元设计两输入 AND、OR 和一输入 NOT，给出连接性质与阈值，并核对真值表。
\end{enumerate}}
\route{圈“M-P”“抑制”“veto”“普通加权”“达到阈值”。查 N01，先判题目是哪种负连接语义，再算；不能将二者混成一种规则。}
\given{M-P 中 active inhibitory input 强制输出0；仅在无抑制输入激活时才比较兴奋输入个数与 \(\theta\)。普通 TLU 使用 \(s=\bm w\T\bm x\)，输出 \(H(s-\theta)\)，没有额外否决。}
\solve
\begin{enumerate}
\item M-P：\(x_3=1\) 是激活的抑制输入，故立即 \(y=0\)，不管另两个兴奋输入有多强。普通 TLU：\(s=1+1-1=1\ge\theta\)，故 \(y=1\)。负权贡献能被正贡献抵消，veto 不能。
\item AND：两输入均兴奋、\(\theta=2\)；OR：两输入均兴奋、\(\theta=1\)。它们的兴奋和均为 \(x_1+x_2\)，只因阈值不同产生不同功能。
\end{enumerate}
\begin{center}\begin{tabular}{ccccc}\toprule
\(x_1\)&\(x_2\)&兴奋和&AND \(\theta=2\)&OR \(\theta=1\)\\\midrule
0&0&0&0&0\\0&1&1&0&1\\1&0&1&0&1\\1&1&2&1&1\\\bottomrule
\end{tabular}\end{center}
NOT 取一个抑制输入、\(\theta=0\)：\(x=1\) 时 veto 输出0；\(x=0\) 时没有抑制，兴奋和为0，且 \(0\ge0\)，所以输出1。普通 TLU 也可用 \(w=-1,\theta=0\) 实现这个特定 NOT，但这不使所有普通负权变成 veto。
\begin{center}\begin{tabular}{ccc}\toprule
\(x\)&是否 veto&NOT输出\\\midrule 0&否，\(0\ge0\)&1\\1&是&0\\\bottomrule
\end{tabular}\end{center}
\ans{给定输入下 M-P 输出0，普通 TLU 输出1。AND：两兴奋输入、阈值2；OR：两兴奋输入、阈值1；NOT：一个抑制输入、阈值0。所有边界按“\(\ge\)”触发，且 veto 优先。}
\qcheck{对每个门核对全部输入而非只试一行。特别检查 OR 的和1、AND 的和2和 NOT 的无输入兴奋和0，均恰在边界，应按约定输出1。}
\pitfall{先求带 \(-1\) 的总和再处理 M-P，导致 veto 被抵消；把“严格大于”误当课件“不小于”；觉得无兴奋输入就绝不能输出1而算错 NOT。}
\transfer{\(n\) 输入 AND 取 \(n\) 个兴奋输入与 \(\theta=n\)，OR 仍为 \(\theta=1\)。若普通 TLU 的权重与阈值同乘正数，预测不变；负数缩放会改变不等号方向，须重新判断。}

\problem{Q08}{Decoder：从目标模式设计并核对全部输入}{正面第2栏 N01；新课件实际第57--61页}
\qtext{构造一个三输入 decoder，仅在 \(\bm r=(1,0,1)\) 时输出1。分别给出 M-P veto 实现和普通加权 TLU 实现；列出完整八行真值表。说明普通 TLU 对一般目标 \(\bm r\in\{0,1\}^n\) 的构造，以及为什么只识别一个模式。}
\route{圈“仅在”“完整模式”“decoder”。查 N01；本题规则已知，做参数构造与真值表验证，不需要感知机训练。先区分 veto 与可相抵的普通负权。}
\given{输入均0/1，输出规则达到阈值触发。目标1的位为第1、3位，目标0的位为第2位；普通 TLU 的分数记 \(s\)，避免与激活后的0/1输出混淆。}
\solve
M-P：第1、3位兴奋，第2位抑制，\(\theta=2\)。若 \(x_2=1\) 立即否决；若 \(x_2=0\)，仍须 \(x_1+x_3\ge2\)，所以必须 \(x_1=x_3=1\)。这正好是指定模式。

普通 TLU：取 \(\bm w=(1,-1,1)\T\)、\(\theta=2\)，即 \(y=H(x_1-x_2+x_3-2)\)。完整核对如下；两种构造的最终输出相同。
\begin{center}\begin{tabular}{ccccc}\toprule
\(x_1\)&\(x_2\)&\(x_3\)&普通分数 \(s=x_1-x_2+x_3\)&两实现输出\\\midrule
0&0&0&0&0\\0&0&1&1&0\\0&1&0&$-1$&0\\0&1&1&0&0\\
1&0&0&1&0\\1&0&1&2&1\\1&1&0&0&0\\1&1&1&1&0\\\bottomrule
\end{tabular}\end{center}
一般构造令 \(k=\sum_i r_i\) 为目标中1的个数，取
\[
w_i=2r_i-1,\qquad \theta=k,\qquad y=H(\bm w\T\bm x-k).
\]
完全匹配时，目标1的位各贡献1，目标0的位贡献0，故总分 \(k\)。目标1误成0，贡献少1；目标0误成1，贡献变成 \(-1\)，也少1。因此每错一位总分就从 \(k\) 减1，只要有错位便低于阈值。这是该简单构造的教学推导。
\ans{M-P：\(x_1,x_3\) 兴奋、\(x_2\) 抑制、\(\theta=2\)。普通 TLU：\(\bm w=(1,-1,1)\T,\theta=2\)。两者均只在101输出1。一般目标可取 \(w_i=2r_i-1,\theta=\sum_i r_i\)，任一错位使分数降低。}
\qcheck{至少检查目标101、含激活抑制的111、少一个所需兴奋的001；完整表确认没有遗漏其他触发输入。参数不必唯一，真正验收条件是全部行正确。}
\pitfall{目标为0的位直接不连接，会把“该位必须0”误成“不关心该位”；把单次匹配识别理解成某种学习收敛；沿用旧课件页码找错 decoder。}
\transfer{若目标改成001，普通参数为 \((-1,-1,1)\T,\theta=1\)；若全零目标，则 \(k=0\)，全负权、阈值0恰好只接收全零。多个正模式可各设 decoder，再用 OR 合并。}

\problem{Q09}{多层 XOR：表示能力与单元限制}{正面第2栏 N01/P02；正面第3栏 B01}
\qtext{输入 \((x_1,x_2)\in\{0,1\}^2\)，XOR 在01、10输出1，在00、11输出0。用两层 M-P 计算单元构造 XOR，写出隐藏量及真值表。再证明任意普通单个加权阈值单元也不能实现 XOR。最后说明“可以构造”是否意味着“某训练算法必能学到且泛化良好”。}
\route{圈“XOR”“两层”“单个阈值”“不能”。查 N01 的正行 decoder→OR，P02 的线性可分限制及 B01 的隐藏层表示；不要反复调单个单元的学习率来寻找本不存在的解。}
\given{所有激活取达到阈值输出1，输入层只传值，不算计算层。第一层设两个模式 detector \(h_1,h_2\)，第二层输出 \(y\)。单个 TLU 使用任意实数 \(w_1,w_2,\theta\)，并非只允许正权。}
\solve
取
\[
h_1=x_1\land\neg x_2,\qquad
h_2=\neg x_1\land x_2,\qquad y=h_1\lor h_2.
\]
\(h_1\)：\(x_1\) 兴奋、\(x_2\) 抑制，\(\theta=1\)；\(h_2\) 对调兴奋与抑制，\(\theta=1\)；输出单元两个兴奋输入、\(\theta=1\)。因此总共三个计算单元、两层计算，隐藏层把两个不同模式分别识别。
\begin{center}\begin{tabular}{ccccc}\toprule
\(x_1\)&\(x_2\)&\(h_1\)&\(h_2\)&\(y=h_1\lor h_2\)\\\midrule
0&0&0&0&0\\0&1&0&1&1\\1&0&1&0&1\\1&1&0&0&0\\\bottomrule
\end{tabular}\end{center}
若单个普通 TLU 可以实现它，四行要求
\[
0<\theta,\quad w_2\ge\theta,\quad w_1\ge\theta,
\quad w_1+w_2<\theta.
\]
中间两式相加给 \(w_1+w_2\ge2\theta>\theta\)，与最后一式矛盾。因此偏置 \(b=-\theta\) 或允许负权也不能解除单个线性阈值单元的限制。多层构造增加的是表示能力。

此构造给出一套能实现真值表的参数，尚未证明某种优化算法能自动找到它；有限四行表示正确也未提供任意实际数据的泛化保证。表示、优化、泛化三个问题须分别论证。
\ans{用两个带否定的 AND detector 加一个 OR 即实现 XOR；单个 TLU 的四个阈值约束彼此矛盾，所以无法实现。两层网络存在可用参数，不等于任意初始化、训练法或新数据都能成功。}
\qcheck{逐行检查隐藏输出，尤其11时两隐藏单元都因对方抑制而输出0。检查“不可能证明”没有偷偷限制权重非负，且保留负类严格小于与正类不小于的边界区别。}
\pitfall{将两个原输入直接 OR；把两层计算误写成只有两个计算单元；把 BP 当成与 MLP 并列的模型；将单层 XOR 训练不终止归咎于训练轮数不足。}
\transfer{任意 Boolean 函数都可为每个正行设 decoder 后 OR；这种直接构造可能很大，并不保证最少单元。无抑制兴奋网络只产生单调输出，不能构造 NOT 这样的下降关系。}

\problem{Q10}{循环网络：用旧状态作同步更新}{正面第2栏 N02/N01；新课件实际第62--64页}
\qtext{本题规定同步离散更新：每个单元延迟一步，所有新值均用旧状态。初态 \((A_0,B_0)=(0,0)\)。输入序列为 \(u_0,u_1,u_2,u_3=1,0,1,0\)，状态规则给定为
\[
A_{t+1}=(u_t\lor A_t)\land\neg B_t,
\qquad B_{t+1}=(u_t\land A_t)\land\neg B_t.
\]
每次更新后以 \(B_{t+1}\) 为本次输出。列出状态和输出；说明为什么它有历史依赖，以及有限自动机与循环 M-P 网络的关系。}
\route{圈“初态”“同步”“旧状态”“更新后输出”。查 N02 的状态递推步骤，再查 N01 的逻辑门。原课件图没有完整给出时序，本题已补明，不能改成逐个立即使用新状态。}
\given{\(u_t\in\{0,1\}\) 是外部输入；\((A_t,B_t)\) 是内部状态，均二值。每步先冻结 \((A_t,B_t)\)，分别算两个新值，最后同时写回。这里改变的是状态，权重未学习。}
\solve
\begin{center}\begin{tabular}{ccccc}\toprule
\(t\)&\(u_t\)&旧状态 \((A_t,B_t)\)&新状态 \((A_{t+1},B_{t+1})\)&输出\\\midrule
0&1&(0,0)&(1,0)&0\\
1&0&(1,0)&(1,0)&0\\
2&1&(1,0)&(1,1)&1\\
3&0&(1,1)&(0,0)&0\\\bottomrule
\end{tabular}\end{center}
首步 \(A_1=(1\lor0)\land1=1\)，\(B_1=(1\land0)\land1=0\)。第二步外部输入为0，但 \(A\) 的自反馈让它保持1。第三步 \(u_2=A_2=1\) 且 \(B_2=0\)，所以 \(B_3=1\)。第四步旧 \(B_3=1\) 同时否决两个新状态，清零。

相同当前输入 \(u=1\)，旧状态00产生新状态10、输出0；旧状态10则产生11、输出1。输出的差别由历史保存的状态引起，而非当前输入不同。只有输出位 \(B\) 也不足以代表完整状态，因为00与10的 \(B\) 同为0，但对下一输入反应不同。

有限自动机只有有限状态和有限种输入，下一状态由旧状态和当前输入决定。用有限个 Boolean 位编码状态，再以 M-P 网络实现每个下一状态位的 Boolean 逻辑，配合反馈与适当延迟，即能模拟有限自动机。课件的命题是表示／模拟能力，不是权重自动学习的保证。本例理论有4个二位状态，从00按给定规则可达00、10、11。
\ans{新状态依次10、10、11、00，更新后输出为0、0、1、0。反馈保存历史，当前输入与旧状态共同决定下一步；有限自动机可由循环 M-P 网络模拟，但须规定状态编码与时序。}
\qcheck{每一行的旧状态都等于上一行的新状态；计算 \(B_1\) 时使用 \(A_0=0\)，不能偷用 \(A_1=1\)。最后旧 \(B=1\) 的 veto 应使两新状态均0。}
\pitfall{把状态更新当权重训练；忽略初态；只给输入不规定同步/异步；把有限状态误说成可记住任意长历史的全部细节。}
\transfer{若初态改成10，第一次 \(u=1\) 就会输出1，必须重算整条状态链。若连续输入全1，本规则从00按00→10→11→00循环，不能无条件称为对任意脉冲流精确“除以二”。}

\problem{Q11}{M-P V.S. Perceptron：共同点、差异与限制}{正面第2栏 D01/N01/P01/P02}
\qtext{用完整句子比较 M-P 神经元与简单感知机。答案须包含共同点、连接／激活与学习方面的差异、适用任务和限制。评价三句话：(1)“它们都有阈值，所以负权一定有 veto”；(2)“感知机只要训练足够久，就能在原二维输入上学会 XOR”；(3)“M-P 固定参数不学习，因此不能被组织成循环网络”。}
\route{圈“compare”“共同点”“区别”“限制”。查 D01，分别按模型定义、参数来源、任务、能力比较；不能仅写“都是神经网络”，也不能把一个单元的限制直接扩大为所有多层网络。}
\given{本课 M-P 指0/1、兴奋／绝对抑制与固定阈值逻辑的定义；感知机采用实数 \(\bm w,b\)，\(y=H(\bm w\T\bm x+b)\)。若写学习式，先声明标签为0/1、\(H(0)=1\)。}
\solve
两者都把多项输入连接到计算单元，经阈值产生二值输出；都可放入更大的网络，并以逻辑／几何方式分析表示能力。这是比“都是机器学习”更具体的共性；M-P 本身未规定从样本学习。
\begin{center}\begin{tabular}{>{\raggedright\arraybackslash}p{.16\linewidth}>{\raggedright\arraybackslash}p{.34\linewidth}>{\raggedright\arraybackslash}p{.37\linewidth}}\toprule
比较维度&M-P&简单 Perceptron\\\midrule
参数与连接&按逻辑设计固定兴奋、抑制连接及阈值&实数权重和偏置，可根据带标签样本调整\\
负作用&激活的抑制输入直接 veto&一般负权降低净输入，可被其他正贡献抵消\\
学习规则&该单元定义没有给训练法&误差纠正，例如0/1约定下 \(\Delta\bm w=\eta(z-y)\bm x\)，\(\Delta b=\eta(z-y)\)\\
典型任务&已知逻辑门、decoder、有限状态逻辑构造&监督二分类，学习一个线性阈值边界\\
能力限制&单个单元受其连接和阈值形式限制；可多层组合&单个原空间线性阈值边界，XOR不可分\\\bottomrule
\end{tabular}\end{center}
感知机收敛保证需要有限训练集线性可分及相应更新条件；不保证任意数据都收敛、最大间隔或未见样本全正确。M-P 的多层构造能扩展 Boolean 表示能力，但不等于自带训练算法。

(1)错：阈值激活与 veto 是两个不同规则，普通负权没有一票否决语义。(2)错：XOR 对单个线性阈值单元不可表示，增加轮数不产生可行参数。(3)错：连接结构与是否训练是不同维度；固定权重的循环网络仍能更新内部状态。
\ans{M-P与感知机共同使用连接、聚合和二值阈值输出；前者强调固定逻辑构造和绝对抑制，后者以可学习实数权重／偏置完成监督线性分类。单个感知机不能表示XOR；多层表示、参数学习与循环状态是不同设计维度。三句话均不正确。}
\qcheck{答案是否同时覆盖共性、明确维度的差异、任务与条件；是否将“不能学会XOR”限定为原输入上的单个线性阈值单元；是否把收敛与泛化分开。}
\pitfall{称 M-P 与感知机都必须从数据训练；把普通负权解释成 veto；称感知机训练求唯一解或最大间隔；以为“循环”必然意味着“在线学习”。}
\transfer{若题给普通加权 TLU 也称“M-P”，先按题目明确公式作答，并指出与课件 veto 定义的不同。改成多层网络后，必须重新讨论表示能力，不能沿用单单元XOR结论。}

\problem{Q12}{混合纠错：由文字建立回归式、梯度与解集}{正面第1栏 F02/M01/M02/M03；梯度学习另见正面第3栏 A01}
\qtext{欲用两个实数特征预测一个连续数值，模型含偏置，按全部样本半平方误差之和评价。给两样本 \(\bm x_1=(1,0)\T,z_1=1\)，\(\bm x_2=(0,1)\T,z_2=2\)。一份解答声称：“\(f(\bm x)=\bm w\T\bm x+b\) 即使 \(b\ne0\) 仍是线性变换；损失梯度为 \(\sum_i(z_i-f_i)\bm x_i\)，按负梯度更新；梯度零表示最大值；两条方程能唯一确定 \(w_1,w_2,b\)。”请逐项纠正，写出模型与目标，从零参数、\(\eta=0.1\) 作一次批量梯度更新，并求全部零误差解。}
\route{圈“连续数值”“含偏置”“全部样本半平方误差”“批量”“唯一”。查 M01/M02/M03 与 F02；本题是可微实数回归，不能把最终输出换成感知机类别后再算损失。}
\given{未知为 \(\bm w=(w_1,w_2)\T\)、\(b\)。预测 \(f_i=\bm w\T\bm x_i+b\)，目标为 \(z_i\)。本题用损失\emph{和}，没有再除样本数；批量梯度中的所有残差均由同一套旧参数产生。}
\solve
由文字建立
\[
f(\bm x)=\bm w\T\bm x+b,\qquad
\min_{\bm w,b}J=\tfrac12\sum_{i=1}^2(f_i-z_i)^2
=\tfrac12(w_1+b-1)^2+\tfrac12(w_2+b-2)^2.
\]
有非零偏置时为仿射映射。求导得到
\[
\nabla_{\bm w}J=\sum_i(f_i-z_i)\bm x_i,
\qquad \frac{\partial J}{\partial b}=\sum_i(f_i-z_i).
\]
原解答把梯度与负梯度符号混了；若用 \(z_i-f_i\)，应将其作为下降\emph{更新方向}相加，而不能再取一次负号。

旧参数全零时 \(f_1=f_2=0\)，所以
\[
\nabla_{\bm w}J=(-1,-2)\T,\quad J_b=-3,
\quad \bm w^+=(0.1,0.2)\T,\quad b^+=0.3.
\]
回算新预测 \(f_1^+=0.4,f_2^+=0.5\)，新损失
\[
J^+=\tfrac12[(0.4-1)^2+(0.5-2)^2]
=\tfrac12(0.36+2.25)=1.305,
\]
低于旧损失 \(J=\tfrac12(1+4)=2.5\)。这是给定步长的一次下降检查，不是对任意步长的保证。

全部零误差解满足
\[
\left[\begin{array}{ccc|c}1&0&1&1\\0&1&1&2\end{array}\right],
\qquad w_1+b=1,\quad w_2+b=2.
\]
系数秩2、未知数3，无矛盾但有自由变量。令 \(b=t\)，得 \(w_1=1-t,w_2=2-t\)，\(t\in\R\)。这些解均令非负平方和为0，故均为全局最小点，说明本题最小点并不唯一；“驻点就是最大”也错误。
\ans{正确梯度是预测减目标。一次更新为 \((w_1,w_2,b)=(0.1,0.2,0.3)\)，新损失1.305。全部零误差最小解为 \((1-t,2-t,t)\)，\(t\in\R\)。非零偏置使映射为仿射；驻点、极值性质与唯一性须分别判断。}
\qcheck{重算两新预测和损失；把任意 \(t\) 的解代回两样本，均得1、2。梯度的权重部分是二维，偏置部分是标量；每个样本同时影响偏置，因此偏置梯度是两个残差之和。}
\pitfall{用新 \(w_1\) 计算第二个样本梯度，意外改成在线更新；将损失和换成平均却沿用同一数值梯度；因有零误差解就推断未来样本预测唯一。}
\transfer{如加正则化，目标和驻点都须重写，不能直接沿用该自由解族；如补第三条独立样本约束，可能得到唯一解，也可能冲突，仍须消元检查。}
